% This work is made available under the terms of the 
% Creative Commons Attribution-ShareAlike 4.0 license,
% http://creativecommons.org/licenses/by-sa/4.0/.

\documentclass[a4paper]{book}

\usepackage{wrapfig}
\usepackage{graphicx}
\usepackage{hyperref}
\usepackage{multirow}
\usepackage{scalefnt}
\usepackage{tikz}

% watermark -- for draft stage
%\usepackage[firstpage]{draftwatermark}
%\SetWatermarkLightness{0.9}
%\SetWatermarkScale{5}

% This work is made available under the terms of the
% Creative Commons Attribution-ShareAlike 4.0 license,
% http://creativecommons.org/licenses/by-sa/4.0/.
%
% Version: $Revision: 2916 $

\newenvironment{tight_itemize}{
\begin{itemize}
  \setlength{\itemsep}{1pt}
  \setlength{\parskip}{0pt}
  \setlength{\parsep}{0pt}}{\end{itemize}
}

\newenvironment{tight_enumerate}{
\begin{enumerate}
  \setlength{\itemsep}{1pt}
  \setlength{\parskip}{0pt}
  \setlength{\parsep}{0pt}}{\end{enumerate}
}

% if you just need a simple heading
% Usage:
%   \heading{the text of the heading}
\newcommand{\heading}[1]{
  \vspace{0.3cm} \noindent \textbf{#1} \newline
}

\newcommand{\icon}[1]{\tikz[baseline=-3pt]\node[inner sep=0pt,outer sep=0pt]{\includegraphics[height=1.1em]{#1}};}


\title{
  \textbf{ADAMS} \\
  {\Large \textbf{A}dvanced \textbf{D}ata mining \textbf{A}nd \textbf{M}achine
  learning \textbf{S}ystem} \\
  {\Large Module: adams-imaging-djl} \\
  \vspace{1cm}
  \includegraphics[width=2cm]{images/imaging-djl-module.png} \\
}
\author{
  Peter Reutemann
}

\setcounter{secnumdepth}{3}
\setcounter{tocdepth}{3}

\begin{document}

\begin{titlepage}
\maketitle

\thispagestyle{empty}
\center
\begin{table}[b]
	\begin{tabular}{c l l}
		\parbox[c][2cm]{2cm}{\copyright 2025} &
		\parbox[c][2cm]{5cm}{\includegraphics[width=5cm]{images/coat_of_arms.pdf}} \\
	\end{tabular}
	\includegraphics[width=12cm]{images/cc.png} \\
\end{table}

\end{titlepage}

\tableofcontents
%\listoffigures
%\listoftables

%%%%%%%%%%%%%%%%%%%%%%%%%%%%%%%%%%%
\chapter{Tools}
This module extends the imaging support using the Deep Java Library\cite{djl}, an open source
library that makes Deep Learning available in Java.

%%%%%%%%%%%%%%%%%%%%%%%%%%%%%%%%%%%
\chapter{Image segmentation}
\section{SAM2}
SAM2\cite{sam2} (Segment Anything Model 2) is a very useful tool for aiding human annotators in
annotating images. The tool generates object masks using one or more points as prompt
within the object that is to be annotated.

The \textit{ImageSegmentationAnnotator} transformer offers a SAM2
tool in their toolboxes. In order to use the tool, you need to do the following:

\begin{tight_itemize}
  \item In the annotation interface, select the \textit{SAM2} tool (the icon shows \textit{SA2}).
  \item Adjust any parameters if necessary, such as model name (\textit{tiny} is recommended on a CPU-only machine)
  and minimum probabilities.
  \item With the left mouse button, perform one or more left-clicks to choose the prompt points (use the left mouse
  button while holding CTRL to remove the current points and start over) and hit \textit{ENTER} to have this image analyzed.
  \item If successful, the annotation(s) will be added to the current ones.
\end{tight_itemize}


%%%%%%%%%%%%%%%%%%%%%%%%%%%%%%%%%%%
% This work is made available under the terms of the 
% Creative Commons Attribution-ShareAlike 4.0 license,
% http://creativecommons.org/licenses/by-sa/4.0/.

\begin{thebibliography}{999}
	% to make the bibliography appear in the TOC
	\addcontentsline{toc}{chapter}{Bibliography}

    % references
	\bibitem{adams}
		\textit{ADAMS} -- Advanced Data mining and Machine learning System \\
		\url{https://adams.cms.waikato.ac.nz/}{}

	\bibitem{djl}
		\textit{Deep Java Library} -- open source library to build and deploy deep learning in Java. \\
		\url{https://djl.ai/}{}

\end{thebibliography}


\end{document}
